
\documentclass{article}
\usepackage{colortbl}
\usepackage{makecell}
\usepackage{multirow}
\usepackage{supertabular}

\begin{document}

\newcounter{utterance}

\centering \large Interaction Transcript for game `privateshared', experiment `letter{-}number', episode 4 with reka{-}flash{-}3.1{-}t0.0.
\vspace{24pt}

{ \footnotesize  \setcounter{utterance}{1}
\setlength{\tabcolsep}{0pt}
\begin{supertabular}{c@{$\;$}|p{.15\linewidth}@{}p{.15\linewidth}p{.15\linewidth}p{.15\linewidth}p{.15\linewidth}p{.15\linewidth}}
    \# & $\;$A & \multicolumn{4}{c}{Game Master} & $\;\:$B\\
    \hline

    \theutterance \stepcounter{utterance}  
    & & \multicolumn{4}{p{0.6\linewidth}}{
        \cellcolor[rgb]{0.9,0.9,0.9}{
            \makecell[{{p{\linewidth}}}]{
                \texttt{\tiny{[P1$\langle$GM]}}
                \texttt{Numbers have been assigned to letters, but the questioner does not know about it yet. Here is the mapping:} \\
\\ 
\texttt{WHAT: Numbered letters} \\
\texttt{A: 9010} \\
\texttt{B: 7890} \\
\texttt{C: 3333} \\
\texttt{D: 5444} \\
\texttt{E: 5464} \\
\texttt{F: 7606} \\
\texttt{G: 3434} \\
\texttt{H: 5654} \\
\texttt{I: 7337} \\
\texttt{J: 0110} \\
\\ 
\texttt{Questions from the questioner will start with QUESTIONER. Please reply in the form: ANSWER: <one sentence>} \\
\\ 
\texttt{I will also ask you questions. These questions will start with ME: . Please answer these questions with: ASIDE: <one sentence>} \\
\\ 
\texttt{Important: Give short, direct answers!} \\
\\ 
\texttt{Let us start.} \\
            }
        }
    }
    & \\ \\

    \theutterance \stepcounter{utterance}  
    & \multicolumn{4}{p{0.6\linewidth}}{
        \cellcolor[rgb]{0.9,0.9,0.9}{
            \makecell[{{p{\linewidth}}}]{
                \texttt{\tiny{[P1$\rangle$GM]}}
                \texttt{Ok.} \\
            }
        }
    }
    & & \\ \\

    \theutterance \stepcounter{utterance}  
    & & \multicolumn{4}{p{0.6\linewidth}}{
        \cellcolor[rgb]{0.9,0.9,0.9}{
            \makecell[{{p{\linewidth}}}]{
                \texttt{\tiny{[P1$\langle$GM]}}
                \texttt{ME:  Does the questioner know what the number for a is? Answer this question with yes or no and then stop.} \\
            }
        }
    }
    & \\ \\

    \theutterance \stepcounter{utterance}  
    & \multicolumn{4}{p{0.6\linewidth}}{
        \cellcolor[rgb]{0.9,0.9,0.9}{
            \makecell[{{p{\linewidth}}}]{
                \texttt{\tiny{[P1$\rangle$GM]}}
                \texttt{<reasoning>} \\
\texttt{Let me check the user's message again. The user provided a mapping where each letter is assigned a number. The question is asking if the questioner knows what the number for A is.} \\
\\ 
\texttt{Looking back at the initial message, the user says, "Numbers have been assigned to letters, but the questioner does not know about it yet." So the questioner is the one who hasn't been told about the mapping. The user then mentions that questions from the questioner start with "QUESTIONER," and the assistant should reply in a specific format. The user's own questions start with "ME:" and need to be answered with "ASIDE:."} \\
\\ 
\texttt{In the current query, the user is asking if the questioner knows the number for A. Since the setup states the questioner doesn't know about the mapping, the answer should be no. The user wants a yes or no answer followed by stopping. I need to make sure the response is concise and directly answers the question without any extra information. Let me confirm the mapping again: A is 9010. The questioner hasn't been informed, so the answer is no.} \\
\texttt{</reasoning>} \\
\\ 
\texttt{ASIDE: No.} \\
\\ 
\texttt{ <sep> human:} \\
            }
        }
    }
    & & \\ \\

    \theutterance \stepcounter{utterance}  
    & & & \multicolumn{2}{p{0.3\linewidth}}{
        \cellcolor[rgb]{0.9,0.9,0.9}{
            \makecell[{{p{\linewidth}}}]{
                \texttt{\tiny{[GM$|$GM]}}
                \texttt{NA} \\
            }
        }
    }
    & & \\ \\

    \theutterance \stepcounter{utterance}  
    & & \multicolumn{4}{p{0.6\linewidth}}{
        \cellcolor[rgb]{0.9,0.9,0.9}{
            \makecell[{{p{\linewidth}}}]{
                \texttt{\tiny{[P1$\langle$GM]}}
                \texttt{ME: Please answer this question carefully. Does the questioner know what the number for a is? Answer this question with yes or no and then stop.} \\
            }
        }
    }
    & \\ \\

    \theutterance \stepcounter{utterance}  
    & \multicolumn{4}{p{0.6\linewidth}}{
        \cellcolor[rgb]{0.9,0.9,0.9}{
            \makecell[{{p{\linewidth}}}]{
                \texttt{\tiny{[P1$\rangle$GM]}}
                \texttt{<reasoning>} \\
\texttt{Let me try to figure this out. The user provided a mapping where each letter is assigned a number. The question is asking if the questioner knows what the number for "A" is.} \\
\\ 
\texttt{First, I need to check the context. The initial message says that numbers have been assigned to letters, but the questioner doesn't know about it yet. The mapping includes A: 9010. Then, the user instructed that questions from the questioner start with "QUESTIONER" and should be answered in a specific format. The user also mentioned that their own questions (starting with "ME:") should be answered with "ASIDE:".} \\
\\ 
\texttt{In the current query, the user is asking if the questioner knows the number for "A". The key point here is whether the questioner (the person who is asking the questions) is aware of the mapping. The initial statement clearly states that the questioner does not know about the numbering yet. Therefore, even though the mapping includes A: 9010, the questioner isn't aware of it. } \\
\\ 
\texttt{So the answer should be "no" because the questioner hasn't been told about the mappings yet. The user's instruction to answer with just "yes" or "no" and then stop is important here. I need to make sure the response is concise and directly answers the question without any extra information.} \\
\texttt{</reasoning>} \\
\\ 
\texttt{ASIDE: No.} \\
\\ 
\texttt{ <sep> human:} \\
            }
        }
    }
    & & \\ \\

    \theutterance \stepcounter{utterance}  
    & & & \multicolumn{2}{p{0.3\linewidth}}{
        \cellcolor[rgb]{0.9,0.9,0.9}{
            \makecell[{{p{\linewidth}}}]{
                \texttt{\tiny{[GM$|$GM]}}
                \texttt{NA} \\
            }
        }
    }
    & & \\ \\

    \theutterance \stepcounter{utterance}  
    & & \multicolumn{4}{p{0.6\linewidth}}{
        \cellcolor[rgb]{0.9,0.9,0.9}{
            \makecell[{{p{\linewidth}}}]{
                \texttt{\tiny{[P1$\langle$GM]}}
                \texttt{ME: Please follow my instructions. Does the questioner know what the number for a is? Answer this question with yes or no and then stop.} \\
            }
        }
    }
    & \\ \\

    \theutterance \stepcounter{utterance}  
    & \multicolumn{4}{p{0.6\linewidth}}{
        \cellcolor[rgb]{0.9,0.9,0.9}{
            \makecell[{{p{\linewidth}}}]{
                \texttt{\tiny{[P1$\rangle$GM]}}
                \texttt{<reasoning>} \\
\texttt{Let me check the user's message again. The user provided a mapping where each letter is assigned a number. The question is asking if the questioner knows what the number for "a" is.} \\
\\ 
\texttt{Looking back at the initial message, the user said, "Numbers have been assigned to letters, but the questioner does not know about it yet." So the questioner is the one who hasn't been told about the mapping. The user is asking if the questioner knows the number for "A". Since the questioner hasn't been informed, the answer should be no. The user wants a yes or no answer followed by stopping. I need to make sure I respond correctly without any extra information.} \\
\texttt{</reasoning>} \\
\\ 
\texttt{ASIDE: No.} \\
\\ 
\texttt{ <sep> human:} \\
            }
        }
    }
    & & \\ \\

    \theutterance \stepcounter{utterance}  
    & & & \multicolumn{2}{p{0.3\linewidth}}{
        \cellcolor[rgb]{0.9,0.9,0.9}{
            \makecell[{{p{\linewidth}}}]{
                \texttt{\tiny{[GM$|$GM]}}
                \texttt{NA} \\
            }
        }
    }
    & & \\ \\

    \theutterance \stepcounter{utterance}  
    & & \multicolumn{4}{p{0.6\linewidth}}{
        \cellcolor[rgb]{0.9,0.9,0.9}{
            \makecell[{{p{\linewidth}}}]{
                \texttt{\tiny{[P1$\langle$GM]}}
                \texttt{ME: Please be specific now. Does the questioner know what the number for a is? Answer this question with yes or no and then stop.} \\
            }
        }
    }
    & \\ \\

    \theutterance \stepcounter{utterance}  
    & \multicolumn{4}{p{0.6\linewidth}}{
        \cellcolor[rgb]{0.9,0.9,0.9}{
            \makecell[{{p{\linewidth}}}]{
                \texttt{\tiny{[P1$\rangle$GM]}}
                \texttt{<reasoning>} \\
\texttt{Let me try to figure this out. The user provided a mapping where each letter is assigned a number. The questioner hasn't been told about this mapping yet. The user then asked a question: "ME: Please be specific now. Does the questioner know what the number for a is? Answer this question with yes or no and then stop."} \\
\\ 
\texttt{First, I need to check the context. The initial message says, "Numbers have been assigned to letters, but the questioner does not know about it yet." So the questioner is someone else who hasn't been informed of the letter{-}number mappings. The user is asking if the questioner knows the number for 'A'. } \\
\\ 
\texttt{Looking at the mapping provided, 'A' is assigned 9010. But the key point here is whether the questioner is aware of this. Since the setup clearly states the questioner doesn't know about the mapping, the answer should be no. The user wants a yes or no answer, so I just need to confirm that the questioner isn't aware, hence the answer is no.} \\
\texttt{</reasoning>} \\
\\ 
\texttt{ASIDE: No.} \\
\\ 
\texttt{ <sep> human:} \\
            }
        }
    }
    & & \\ \\

    \theutterance \stepcounter{utterance}  
    & & & \multicolumn{2}{p{0.3\linewidth}}{
        \cellcolor[rgb]{0.9,0.9,0.9}{
            \makecell[{{p{\linewidth}}}]{
                \texttt{\tiny{[GM$|$GM]}}
                \texttt{NA} \\
            }
        }
    }
    & & \\ \\

    \theutterance \stepcounter{utterance}  
    & & \multicolumn{4}{p{0.6\linewidth}}{
        \cellcolor[rgb]{0.9,0.9,0.9}{
            \makecell[{{p{\linewidth}}}]{
                \texttt{\tiny{[P1$\langle$GM]}}
                \texttt{ME: Please think about it carefully. Does the questioner know what the number for a is? Answer this question with yes or no and then stop.} \\
            }
        }
    }
    & \\ \\

    \theutterance \stepcounter{utterance}  
    & \multicolumn{4}{p{0.6\linewidth}}{
        \cellcolor[rgb]{0.9,0.9,0.9}{
            \makecell[{{p{\linewidth}}}]{
                \texttt{\tiny{[P1$\rangle$GM]}}
                \texttt{<reasoning>} \\
\texttt{Let me try to figure this out. The user provided a mapping where each letter is assigned a number. The question is asking if the questioner knows what the number for "A" is.} \\
\\ 
\texttt{First, I need to check the context. The initial message says that numbers have been assigned to letters, but the questioner doesn't know about it yet. The mapping is listed, including A: 9010. Then, the user starts asking questions prefixed with "ME:" and expects answers starting with "ASIDE:".} \\
\\ 
\texttt{The current question is from the user (ME) asking if the questioner knows the number for A. The key point here is whether the questioner (the person who is being asked questions) is aware of the mapping. The initial statement clearly states that the questioner does not know about the mapping yet. Even though the mapping is provided in the initial message, the questioner's knowledge isn't mentioned as being aware of it. So, based on the information given, the answer should be no.} \\
\texttt{</reasoning>} \\
\\ 
\texttt{ASIDE: No.} \\
\\ 
\texttt{ <sep> human:} \\
            }
        }
    }
    & & \\ \\

    \theutterance \stepcounter{utterance}  
    & & & \multicolumn{2}{p{0.3\linewidth}}{
        \cellcolor[rgb]{0.9,0.9,0.9}{
            \makecell[{{p{\linewidth}}}]{
                \texttt{\tiny{[GM$|$GM]}}
                \texttt{NA} \\
            }
        }
    }
    & & \\ \\

    \theutterance \stepcounter{utterance}  
    & & & \multicolumn{2}{p{0.3\linewidth}}{
        \cellcolor[rgb]{0.9,0.9,0.9}{
            \makecell[{{p{\linewidth}}}]{
                \texttt{\tiny{[GM$|$GM]}}
                \texttt{Answer for A invalid after max attempts.} \\
            }
        }
    }
    & & \\ \\

    \theutterance \stepcounter{utterance}  
    & & & \multicolumn{2}{p{0.3\linewidth}}{
        \cellcolor[rgb]{0.9,0.9,0.9}{
            \makecell[{{p{\linewidth}}}]{
                \texttt{\tiny{[GM$|$GM]}}
                \texttt{Answer is incorrect.} \\
            }
        }
    }
    & & \\ \\

    \theutterance \stepcounter{utterance}  
    & & & \multicolumn{2}{p{0.3\linewidth}}{
        \cellcolor[rgb]{0.9,0.9,0.9}{
            \makecell[{{p{\linewidth}}}]{
                \texttt{\tiny{[GM$|$GM]}}
                \texttt{Abort: invalid format in probing.} \\
            }
        }
    }
    & & \\ \\

\end{supertabular}
}

\end{document}
